\chapter{Background: Integrity, Canonicalization, Provenance, and Authorization}\label{ch02:background}

This chapter is a tutorial survey of the standard technical material on
which the CAS Core 8 kernel builds: cryptographic hashing and content
addressing, canonical serialization, reproducible builds and
supply-chain integrity, provenance modeling, the
integrity-versus-authorization distinction, threat-modeling method, and
the distinction between process isolation and sandboxing. Everything in
this chapter is established background from the security, systems, and
standards literature. None of it is project novelty, and the chapter's
framing is deliberately impersonal: the kernel appears only in forward
pointers, as a consumer of these ideas, never as their source. The
kernel's specific contribution---discussed in later chapters---is the
composition of standard mechanisms into a fail-closed substrate for a
particular research program, together with a particular design stance
about authorization; the mechanisms themselves are the community's.

\begin{remark}[How to read this chapter]\label{ch02:howtoread}
Readers fluent in applied cryptography and security engineering can
skim: the load-bearing artifacts are the canonicalization comparison
axes of Section~\ref{ch02:canon-axes}, the isolation ladder of
Section~\ref{ch02:isolation}, and the closing map from background
concepts to kernel mechanisms in Section~\ref{ch02:synthesis}. Readers
new to the material should read linearly; each section is written to be
self-contained, defines its terms at first use, and states not only
what each mechanism provides but---because the kernel's design turns on
these negatives---what it conspicuously does not provide.
\end{remark}

\section{Cryptographic hashing and content addressing}\label{ch02:hashing}

\subsection{Hash functions and their contract}\label{ch02:hashcontract}

A cryptographic hash function is a deterministic function $H$ mapping
arbitrary-length byte strings to fixed-length digests, engineered so
that the mapping behaves, for all practical adversarial purposes, like
a random function. Three properties make up the standard contract.
\emph{Preimage resistance}: given a digest $d$, it is computationally
infeasible to find any message $m$ with $H(m) = d$. \emph{Second-preimage
resistance}: given a message $m_1$, it is infeasible to find a distinct
$m_2$ with $H(m_1) = H(m_2)$. \emph{Collision resistance}: it is
infeasible to find any pair $m_1 \neq m_2$ with $H(m_1) = H(m_2)$.
Collision resistance is the strongest and most fragile of the three:
by the birthday bound, a digest of $n$ bits offers at most on the order
of $2^{n/2}$ work against collisions, and historical functions (MD5,
SHA-1) fell to structural attacks far below even that bound. Current
practice therefore standardizes on functions with comfortable margins;
SHA-256, the function used throughout the kernel's sources, produces
256-bit digests, giving a nominal $2^{128}$ birthday bound and, at the
time of writing, no publicly known structural attacks that materially
degrade it.

Alongside the adversarial contract sits a humbler property that does
much of the practical work: the \emph{avalanche} behavior by which any
single-bit change to the input flips, in expectation, half the digest
bits. Avalanche is what makes a digest useful as a fingerprint: near
misses do not look near. Combined with determinism, this yields the
basic evidence pattern of every integrity system: if two computations
of $H$ over allegedly identical bytes disagree, the bytes differ,
full stop; if they agree, the bytes are identical unless the
adversary has broken the function.

Two consequences of the contract deserve emphasis because they are
routinely elided in casual use. First, a hash function is defined over
\emph{bytes}, not over meanings. $H$ knows nothing about JSON, about
Unicode, about numbers, about the intent of the producer. Two encodings
of the same abstract value hash differently unless the encodings are
byte-identical---which is the entire subject of
Section~\ref{ch02:canonical}. Second, a digest carries no information
about who computed it, why, or with what authority. Anyone can hash
anything. A digest is a statement about content, never a statement
about endorsement; the machinery that turns content statements into
origin statements (signatures, authenticated channels, trust roots)
is separate machinery, surveyed in Section~\ref{ch02:integrity-auth},
and conflating the two is one of the classic errors of applied
security engineering \cite{anderson}.

\subsection{Function lifecycle and digest agility}\label{ch02:agility}

Hash functions age. MD5 went from standard to broken-for-collisions
within a decade of wide deployment; SHA-1 followed, first with
theoretical attacks and then with demonstrated collisions; and
systems that had woven either function into long-lived identifiers
faced expensive migrations because a digest, once used as a name, is
referenced everywhere and controlled nowhere. The standard
engineering responses are worth recording as background because any
system that persists digests inherits the problem. First,
\emph{margin}: choose a function whose security comfortably exceeds
present need, since migration cost dwarfs the marginal cost of a
longer digest. Second, \emph{explicitness}: record which function
produced each digest---in a field, a prefix, or a fixed system-wide
declaration---so that a future migration can distinguish old names
from new mechanically rather than by guesswork. Third,
\emph{planned coexistence}: migrations happen by dual-hashing
during a transition window, and schema designs that permit a second
digest field beside the first migrate gracefully, while designs that
hard-wire ``the'' digest do not. Fourth, a caution against
overreach: \emph{agility} in the protocol-negotiation sense (letting
peers choose among functions at run time) has its own failure
modes---downgrade attacks, cross-function confusion where the same
content has multiple live names---and evidence systems generally
prefer a single fixed function per era with explicit migration
events over continuous negotiation. The practical upshot for
digest-bearing records: fix one function, name it, validate digest
syntax strictly (length, alphabet, case convention) at every intake
boundary, and treat any future function change as a schema-visible
event rather than a silent substitution.

\subsection{Digests as names: content addressing}\label{ch02:contentaddr}

Once digests are cheap and collision-resistant, a powerful naming
discipline becomes available: use $H(\text{content})$ \emph{as the
name of the content}. A store keyed this way is called
content-addressed. The discipline has properties that location-based
naming (file paths, URLs, database keys) cannot offer:

\begin{itemize}
\item \emph{Self-verification.} Given a name and alleged content,
anyone can check the binding by rehashing. Retrieval integrity needs
no trusted server: a corrupt or malicious store can fail to produce
the content, but cannot produce the wrong content undetected.
\item \emph{Immutability of reference.} A content address names one
byte string forever. There is no ``the file changed under the same
name''; change produces a new name. Version history becomes a
first-class object rather than an overwrite.
\item \emph{Deduplication and convergence.} Identical content produced
independently receives the same name, so replicas can be reconciled
and duplicates collapsed without coordination.
\item \emph{Composability via hash trees.} Structures whose nodes name
their children by digest---Merkle trees and directed acyclic graphs---
extend the discipline to aggregates: the root digest commits to the
entire reachable structure, and any tampering anywhere invalidates
the root. Version-control systems, transparency logs, and
package-distribution systems are all elaborations of this one idea.
\end{itemize}

\subsection{Hash trees, logs, and the mechanics of aggregation}\label{ch02:merkle}

The hash-tree idea deserves a slightly longer treatment, because the
kernel's registry and receipt structures are degenerate cases of it
and the general form explains the design pressure. In a Merkle tree,
each leaf is the digest of a data block and each internal node is the
digest of its children's digests; the root digest therefore commits
to every leaf, and membership of any leaf can be proven with a
logarithmic-size path of sibling digests. Three practical systems
elaborate the idea in three directions. Version-control systems make
the tree a DAG of commits, each naming its parents and its content
tree by digest, so that a single commit identifier fixes an entire
history; the append-only, tamper-evident character of such histories
is a direct corollary of content addressing. Transparency logs make
the tree an append-only sequence with signed tree heads, so that a
log operator can prove both membership and append-only consistency
to auditors, and equivocation---showing different histories to
different observers---becomes detectable by gossip between observers.
Package ecosystems make the tree a manifest: a release names its
files by digest, a lockfile names its dependencies by digest, and
installation verifies the closure.

Two mechanical details recur in all three and matter for any system
that records digests in evidence. First, \emph{domain separation}:
when digests of digests are computed, leaf and node encodings must be
distinguishable (by a prefix byte or a typed encoding), or an
attacker can present an internal node as a leaf and forge
memberships; the general lesson is that even hashing needs typing.
Second, \emph{presentation discipline}: a 256-bit digest has many
textual presentations---upper and lower hexadecimal, base64, prefixed
and truncated forms---and systems that compare digests as strings
must fix one presentation or normalize before comparison, since
case-mismatched hexadecimal is byte-different while
digest-identical. These are small points, but evidence systems live
and die by small points; a registry that stores lowercase digests
and receives uppercase ones will either fail closed (annoying,
recoverable) or compare case-insensitively somewhere and exactly
somewhere else (a latent inconsistency of exactly the kind canonical
disciplines exist to forbid). Fixing a single normalized digest
presentation at every interface is the standard remedy.

The limits are the duals of the properties. Content addressing binds
names to bytes, so it inherits byte-sensitivity: a semantically
irrelevant byte change (a rewritten line ending, a reordered map key)
produces a new name and severs every reference. Systems therefore face
a standing temptation to normalize content before hashing---and every
normalization is a decision about which distinctions are ``irrelevant,''
a decision that itself needs to be explicit, deterministic, and
recorded, or the cure reintroduces the disease. The kernel's
two-identity discipline for line endings (raw digest and normalized
comparison digest carried as separate, non-substitutable fields) is a
forward pointer here: it is one project-specific answer to a
completely standard tension.

A second limit: content addressing authenticates content, not
context. Knowing that bytes match a digest says nothing about whether
those bytes are the \emph{right} bytes for a purpose---the current
version, the reviewed version, the version this receipt was actually
computed over. Binding digests to purposes is the job of provenance
(Section~\ref{ch02:prov}) and of typed records that carry digests in
named fields with declared roles.

\section{Canonical serialization}\label{ch02:canonical}

\subsection{The problem}\label{ch02:canon-problem}

Digest-based evidence has a hidden premise: that equal abstract values
produce equal bytes. Serialization formats violate this premise by
default. The same logical map, list, or number admits many legal
encodings, and any process that re-serializes---a different library, a
different version, a different platform---may lawfully produce
different bytes and hence a different digest. The result is the worst
kind of failure for an evidence system: \emph{false negatives}
(identical evidence appearing to differ, breaking replay comparisons
and cache keys) and, more insidiously, \emph{false positives} via type
punning (distinct values encoding to identical bytes, so that
different evidence hashes equal). A canonical serialization is a
function from abstract values to bytes that is total on a declared
domain, deterministic, and injective on that domain: one value, one
byte string; distinct values, distinct byte strings.

\begin{definition}[Canonical encoding]\label{ch02:def-canonical}
Let $V$ be a declared domain of abstract values. A \emph{canonical
encoding} for $V$ is a function $c : V \to B$ into byte strings that
is (i) total on $V$, (ii) deterministic across implementations,
platforms, and runs, and (iii) injective: $c(v_1) = c(v_2)$ implies
$v_1 = v_2$ under the domain's declared equality. A system adopting
$c$ additionally declares its \emph{out-of-domain behavior}: what
happens when a value outside $V$ is presented.
\end{definition}

Under Definition~\ref{ch02:def-canonical}, digest evidence acquires
its intended meaning: for $v_1, v_2 \in V$, $H(c(v_1)) = H(c(v_2))$
holds exactly when $v_1 = v_2$, up to hash collisions. The three
clauses fail independently in practice, and each failure has its own
signature: partiality surfaces as crashes or silent omissions on
exotic values; nondeterminism surfaces as false difference; and
non-injectivity surfaces as false identity, the quietest and most
dangerous of the three. The subsections that follow survey the
standard obstacles to each clause, because each obstacle corresponds
to a real divergence observed across mainstream serializers.

\subsubsection{Key order}\label{ch02:canon-keys}

JSON objects and most map types are semantically unordered, but their
serializations are sequences; some order must be chosen. Libraries
default variously to insertion order, hash-iteration order (which can
vary across runs, processes, and library versions, and in some
languages is deliberately randomized), or sorted order under some
collation. Any digest computed over an encoding with unspecified key
order is nondeterministic in the fleet sense: it may be stable on one
machine today and differ on another tomorrow. A canonical form must
fix a total order on keys and say precisely over what it sorts---bytes
of the UTF-8 encoding, UTF-16 code units, Unicode code points---since
these orders differ once strings leave ASCII. A related trap is
duplicate keys: JSON's grammar permits them, parsers disagree on
whether the first or last wins, and an encoder that tolerates them has
already lost injectivity at the parsing boundary. Strict systems
reject duplicates outright.

\subsubsection{Numbers}\label{ch02:canon-numbers}

Numbers are the richest source of canonicalization failure. Binary
floating point admits multiple decimal renderings of the same value
(\texttt{1E2}, \texttt{100.0}, \texttt{1.0e+2}), and conversely,
distinct values with renderings that careless formatting collapses.
The modern remedy for rendering is the shortest-round-trip rule:
render the fewest digits that parse back to the identical binary
value. But even with rendering fixed, deeper puns remain. Is the
integer $1$ the same value as the float $1.0$? Many serializations
say yes by encoding both as \texttt{1}; many type systems say no. Is
$-0$ the same as $0$? IEEE~754 distinguishes them; most decimal
renderings do not. \texttt{NaN} and the infinities have no JSON
representation at all, so libraries variously reject them, encode
\texttt{null}, or emit non-standard tokens---three behaviors, three
different digests, one input. Exact decimal quantities (the natural
carrier for measured and monetary values) add a further layer: the
decimal value $1.0$, the string \texttt{"1.0"}, and the integer $1$
are three distinct typed values that naive encodings can collapse
into one byte string. An evidence encoding must decide, explicitly,
which numeric distinctions are real, and then either represent them
distinguishably (for example by tagging exact decimals as typed
objects) or exclude the ambiguous carriers from the domain entirely
(for example by rejecting binary floats). Both are defensible; the
indefensible option is the default one, silence.

\subsubsection{Strings and Unicode}\label{ch02:canon-unicode}

Unicode offers multiple byte sequences for perceptually and even
canonically equivalent text: precomposed versus decomposed forms
(NFC/NFD), compatibility variants (NFKC/NFKD), invisible format
controls, bidirectional override characters, and unpaired surrogates
that some decoders pass through and others reject. Escaping rules
multiply encodings further: JSON permits \texttt{"/"} and
\texttt{"\textbackslash/"}, permits any character to be
\texttt{\textbackslash u}-escaped, and permits both short and long
escape forms. A canonical form must pick one escaping policy and one
normalization stance. There is also a security dimension: invisible
and bidirectional characters enable identifiers that render
identically but differ in bytes---a spoofing vector whenever humans
review machine identifiers. The conservative response, standard in
review-critical systems, is to split the string domain: machine
identifiers are restricted to a declared ASCII profile (trimmed,
normalized, drawn from an explicit alphabet), while unrestricted
prose travels in separate fields that never participate in identifier
comparison.

\subsubsection{Aggregates, sharing, and cycles}\label{ch02:canon-cycles}

Serialization formats describe trees; program values are graphs. Two
structural questions follow. First, \emph{unordered containers}: a
set has no traversal order, so any linearization imposes one; unless
the element type carries a total order that the encoder applies, set
encodings are nondeterministic, and strict systems simply reject sets
from the evidence domain. Second, \emph{sharing and cycles}: a value
reachable twice can be encoded twice (losing sharing) or by
reference (requiring a reference scheme), and a cyclic value cannot
be tree-encoded at all---a naive recursive encoder will not
terminate. Robust canonicalizers are therefore explicitly bounded:
they track visited nodes, fail on cycles, and enforce declared depth,
node-count, and output-size limits, converting resource-exhaustion
attacks into clean rejections. Boundedness is not a luxury; an
encoder that can be driven to unbounded recursion or output by a
crafted input is a denial-of-service primitive sitting in the middle
of the evidence path.

\subsubsection{Type identity}\label{ch02:canon-types}

Finally, serialization erases types unless told otherwise. A tuple
and a list, an enumeration member and its underlying value, a
domain-specific wrapper and its payload---each pair collapses to the
same encoding in most formats. Whether that collapse is acceptable
depends entirely on whether the type distinction is load-bearing. In
an evidence system where enumerated statuses are precisely the values
that must not be forgeable by strings, the collapse is not
acceptable, and the standard remedy is type tagging: the encoding of
a typed value embeds a reserved marker naming its type, and the
reserved marker is excluded from ordinary maps so that data cannot
impersonate tags. Tagging trades a little verbosity for injectivity
across type boundaries, and it must be paired with the exclusion
rule---a tagging scheme whose tags can be spoofed from plain data
has merely moved the pun.

\subsubsection{A worked divergence}\label{ch02:canon-worked}

To make the failure concrete, consider one modest abstract value: a
record with a name field, a count, and a measured ratio---say the
name \texttt{"caf\'e"}, the count one, and the ratio one-tenth. A
survey of mainstream serializers, run naively, can lawfully produce
all of the following encodings, each yielding a different SHA-256
digest: keys in insertion order versus sorted order; the name with
the accented character raw (in either NFC or NFD form) versus
escaped (in short or long escape form); the count as \texttt{1}
versus \texttt{1.0} (an integer widened by a numeric tower) versus
\texttt{"1"} (stringified by a lossy channel); and the ratio as
\texttt{0.1} versus \texttt{0.10000000000000001} (a formatter
printing seventeen digits) versus \texttt{1E-1}. That is at least
several dozen byte-distinct encodings of one value, produced by
correct, widely deployed software making defensible local choices.
An evidence pipeline in which any two components disagree on any one
of these choices will report difference where there is identity.
Worse, run the survey in reverse: the string \texttt{"1"} and the
integer \texttt{1}, or a decimal one-tenth and a binary
double one-tenth, can each arrive at the \emph{same} bytes through
some component pair, reporting identity where there is difference.
Every rule in a strict canonical encoding can be read off from this
example: fix the key order, fix the string form and reject the
invisible cases, split the numeric types and reject the ambiguous
carrier, tag what remains, and refuse everything else.

\subsubsection{Lessons from canonicalization history}\label{ch02:canon-history}

The difficulty is also attested historically, and one episode is
instructive enough to summarize as background. The XML signature
stack of the early 2000s required canonicalization (C14N) so that
signatures survived innocuous re-serialization, and the resulting
specifications---canonical XML, exclusive canonicalization, and
their interaction with signature transforms---became notorious for
complexity: namespace inheritance, attribute ordering, whitespace in
mixed content, and processing-instruction handling all demanded
rules, and implementations diverged in the corners. The security
literature subsequently documented signature-wrapping attacks in
which the signed subtree and the processed subtree were different
nodes of the same document---an exploitation not of the hash or the
signature but of the gap between what was canonicalized and what
was consumed. Two standard lessons are drawn from that experience.
First, canonicalization complexity is attack surface: the more
transformations sit between the value the producer meant and the
bytes the verifier hashes, the more room for the two to diverge.
Second, the safest designs canonicalize \emph{once}, at a single
declared boundary, and consume exactly the bytes they verified,
rather than re-deriving the value downstream. JSON-era designs,
including JCS \cite{jcs} and the stricter evidence encoders
surveyed below, are deliberately minimal in reaction: small value
domains, few transformations, one route from value to bytes. The
kernel's one-serialization-route-per-type rule (later chapters) is
this lesson in project form.

\subsection{RFC 8785: the JSON Canonicalization Scheme}\label{ch02:jcs}

The reference point for JSON canonicalization is the JSON
Canonicalization Scheme (JCS), published as RFC~8785 \cite{jcs}. JCS
targets the interoperable subset of JSON (I-JSON): data that fits in
UTF-8, with numbers representable as IEEE~754 double precision and no
duplicate object keys. Within that subset, JCS fixes: object keys
sorted by their UTF-16 code-unit sequences; number serialization per
the ECMAScript shortest-round-trip algorithm; a minimal, fixed string
escaping; no insignificant whitespace; and UTF-8 output. The design
philosophy is deliberate reuse: by matching ECMAScript's own
serialization behavior, JCS lets a large installed base produce
canonical output with modest effort, and it achieves its goal of
byte-stable digests for I-JSON data across conforming
implementations.

Two of the scheme's choices repay a closer look, since they
illustrate how a canonicalization standard turns the obstacles of
the preceding subsections into decisions. The key-ordering rule
sorts object keys by comparing their UTF-16 code-unit sequences.
For ASCII keys this coincides with byte order, but outside the
basic multilingual plane the choice is visible: characters encoded
with surrogate pairs in UTF-16 compare differently under
code-unit order than under code-point or UTF-8 byte order, so the
same object can sort differently under a plausible-but-wrong
implementation and produce a plausible-but-wrong digest. The rule
is a reminder that ``sorted keys'' is not yet a specification;
\emph{sorted by what} is. The number rule adopts the ECMAScript
serialization algorithm, which emits the shortest decimal string
that round-trips to the same IEEE~754 double, with fixed rules for
when exponent notation appears and with negative zero serialized
as \texttt{0}. This buys cross-implementation byte stability for
every representable double, at the declared price of the domain
restriction: values needing more than double precision, and the
non-finite values, are simply outside I-JSON and hence outside the
scheme's promises. Systems whose evidence includes such values must
carry them in strings or tagged structures by their own
convention---which is precisely the pattern the strict encoder
family adopts for exact decimals.

Equally instructive is what JCS deliberately does not do. It does not
extend the value domain: sets, cycles, and shared structure are
outside JSON and hence outside JCS. It does not distinguish numeric
types: within I-JSON, a number is a double, so the integer/decimal/
float distinctions of Section~\ref{ch02:canon-numbers} are collapsed
by the domain choice itself, and values outside double range or
precision are simply not I-JSON. It does not normalize Unicode
content: strings are canonicalized in escaping, not in normalization
form, so NFC and NFD spellings of the same text remain distinct
values. It does not tag types, bound resources, or take a stance on
rejection versus coercion for out-of-domain values---those questions
fall to the systems that adopt it. None of these are defects; they
are scope decisions appropriate to JCS's interoperability goal. But
they mean that JCS is a \emph{component} of an evidence discipline,
not the discipline itself: a system whose evidence includes exact
decimals, typed enumerations, or adversarial inputs must layer domain
restrictions, tagging, and bounds on top of---or instead of---the
interoperability-oriented core.

\subsection{The rejection stance}\label{ch02:canon-rejection}

The out-of-domain axis deserves a short argument of its own, because
it is where evidence encoders and general-purpose serializers part
company most sharply. A general-purpose serializer is judged by how
much it can encode; faced with a set, it sorts or lists it; faced
with a NaN, it emits a null or a token; faced with an unknown
object, it calls a fallback hook. Every one of these accommodations
is a silent semantic decision made on the producer's behalf, and in
an evidence setting each becomes a pun generator: the sorted set is
now indistinguishable from a sorted list; the null-for-NaN is
indistinguishable from a genuine null; the fallback-encoded object
is indistinguishable from a hand-built map that happens to match.
The coercion that makes the serializer convenient makes the digest
ambiguous, and digest ambiguity is unrecoverable downstream---no
consumer of the digest can reconstruct which of the collapsed
values was meant.

The rejection stance inverts the judgment: an evidence encoder is
judged by how little it encodes without instruction. Everything
outside the declared domain fails closed, with an error naming the
offending value and the rule it violates, at the moment and place
where the producer---who alone knows what was meant---can fix it.
The cost profile is the mirror image of coercion's: friction is
paid up front, visibly, per producer, while the benefit (every
digest in the system means exactly one thing) accrues silently,
system-wide, forever. Which stance is right is not absolute but
positional---the same organization may correctly run coercing
serializers at its logging edges and a rejecting encoder at its
evidence core---and the classic failure is stance drift: an
evidence path that, under delivery pressure, imports the convenient
serializer ``temporarily.'' The stance is therefore best enforced
structurally (a single blessed encoder for the evidence domain, no
alternative import path), which is the design shape later chapters
describe for the kernel; the argument for it, as this section has
tried to show, is entirely general.

\subsection{Comparison axes for canonical encodings}\label{ch02:canon-axes}

Surveying the design space, canonical encodings can be compared along
seven axes, summarized in Table~\ref{ch02:tab-axes} (the number-policy
axis occupies two rows there, one for finite and one for non-finite
treatment).

\begin{enumerate}
\item \emph{Value domain}: which abstract values are encodable at
all---full host-language values, a format's native model, or a
deliberately restricted evidence domain.
\item \emph{Number policy}: single numeric type versus distinguished
integers, decimals, and floats; treatment of non-finite values.
\item \emph{String policy}: escaping rules; Unicode normalization
stance; identifier profiles versus free prose.
\item \emph{Ordering}: what is sorted, and over what collation.
\item \emph{Out-of-domain behavior}: coerce (make something
encodable), best-effort (encode with caveats), or reject (fail
closed). This axis, more than any other, separates interoperability
formats from evidence formats: coercion maximizes the chance of
producing output; rejection maximizes the meaning of the output that
is produced.
\item \emph{Type identity}: erased types versus tagged types, and
whether tags are spoofable from data.
\item \emph{Resource bounds}: unbounded recursion versus declared
depth, node, and size limits with cycle detection.
\end{enumerate}

\begin{table}[htbp]
\centering
\caption{Comparison axes for canonical serialization designs. ``JCS''
is RFC~8785 \cite{jcs}; ``strict evidence encoders'' names the family
of reject-and-tag designs of which the kernel's encoder (later
chapters) is one member.}
\label{ch02:tab-axes}
\begin{tabular}{@{}lll@{}}
\toprule
Axis & JCS (RFC 8785) & Strict evidence encoders \\
\midrule
Value domain & I-JSON & declared, restricted \\
Numbers & IEEE double only & int/decimal split; floats rejected \\
Non-finite values & outside I-JSON & rejected explicitly \\
Strings & escape canonicalization & profiles; controls rejected \\
Ordering & UTF-16 code-unit key sort & sorted string keys \\
Out-of-domain & out of scope & fail closed \\
Type identity & erased & reserved-tag typed values \\
Resource bounds & unspecified & depth/node/size bounds; cycle fail \\
\bottomrule
\end{tabular}
\end{table}

The table should be read as a contrast of goals, not of quality. JCS
optimizes for the interoperability of independently implemented
producers over benign data; strict evidence encoders optimize for the
evidentiary meaning of a digest over potentially adversarial data.
The first wants maximal producers; the second wants minimal surprise.
A mature system frequently needs both, in different places---an
interchange layer speaking a JCS-like dialect at its edges and a
strict encoder at its evidentiary core---and confusing the two roles
in either direction (strictness at the interoperability edge,
tolerance at the evidence core) produces respectively friction and
laundering.

\section{Reproducible builds and supply-chain integrity}\label{ch02:repro}

\subsection{The reproducibility contract}\label{ch02:repro-contract}

A build is \emph{reproducible} when, in the Reproducible Builds
project's canonical formulation, ``given the same source code, build
environment and build instructions, any party can recreate bit-by-bit
identical copies of all specified artifacts'' \cite{reprobuilds}. The motivation is
a trust gap at the heart of software distribution: users run
binaries, but review---when it happens at all---happens on sources.
Nothing in a conventional build pipeline connects the two: the
binary's contents are determined not only by the source but by the
compiler, its version and flags, the build machine's state, and
whatever compromised any of them. Reproducibility closes the gap by
making the source-to-binary function deterministic, so that
independent rebuilders can each compile the source and compare
digests. Agreement among independent parties establishes---to the
strength of the hash function and the independence of the
rebuilders---that the distributed binary is what the reviewed source
produces. A single malicious or compromised builder can no longer
substitute a backdoored artifact without every honest rebuilder's
digest disagreeing.

Achieving bit-for-bit determinism is an exercise in hunting down
incidental environment leakage, and the catalog of leaks doubles as
a catalog of what ``environment'' even means: embedded timestamps
and dates; absolute build paths baked into debug info; locale- and
timezone-dependent formatting; filesystem readdir order and archive
member order; parallelism-dependent link order; uninitialized memory
serialized into outputs; randomized identifiers and, notably,
randomized hash-iteration order in the build tooling itself. The
remedies are correspondingly standardized: pinned toolchains,
normalized metadata, forced stable orderings, and declared epochs
for anything time-shaped. The parallel with
Section~\ref{ch02:canonical} is exact and not coincidental:
reproducible builds are canonical serialization applied to the
compilation function, and both disciplines reduce to the same
maxim---every source of nondeterminism is either eliminated,
pinned and recorded, or it silently voids digest comparison.

\subsection{Pinning, lockfiles, and the verification economy}\label{ch02:repro-pinning}

Between the source and the build sits the dependency graph, and the
reproducibility discipline extends to it through pinning: a
dependency manifest that names not version \emph{ranges} but exact
artifacts, each with its digest, so that resolution is a lookup
rather than a negotiation. The lockfile pattern---a generated,
committed file recording the fully resolved, hashed closure of the
dependency graph---is the ecosystem-standard instantiation. Its
evidentiary reading is the important one here: a lockfile converts
``the build used library X'' from a claim about intentions into a
claim about bytes, checkable by rehashing. Unpinned resolution, by
contrast, makes every build a function of the repository state of
third parties at build time, which is exactly the kind of
unrecorded input that voids digest comparison; the supply-chain
attack literature is largely a catalog of ways this openness is
exploited (dependency confusion between internal and public names,
typosquatted packages, maintainer-account compromise followed by a
poisoned patch release), and pinning-plus-hashing is the standard
first line against all of them.

The verification side has its own economics, worth stating because
they explain why reproducibility is a \emph{social} mechanism and
not merely a technical one \cite{reprobuilds}. One reproducible
build verified by nobody provides nothing; the security benefit
scales with the number and independence of rebuilders who actually
recompute and compare, because the attacker must now compromise all
of them---or the comparison channel---rather than one build server.
Independence is doing the work in that sentence: rebuilders sharing
the same compromised toolchain image reproduce each other's
compromise. The design consequences that generalize to any evidence
system are two. Cheap verification is a feature to be engineered
deliberately---the whole pipeline from source to comparable digest
must be runnable by an outsider at tolerable cost, or independence
stays theoretical. And comparison results are themselves evidence
needing the full discipline of this chapter: a rebuild receipt has
content, provenance, and (no) authority, exactly like any other
receipt.

\subsection{Limits and the supply-chain frame}\label{ch02:repro-limits}

Reproducibility, like hashing, has sharply bounded semantics.
Bit-for-bit agreement establishes that binary corresponds to source;
it establishes nothing about whether the source is any good.
Reproducibly built malware is still malware, reproducibly built bugs
are still bugs, and a reproducible pipeline whose declared toolchain
is itself compromised reproduces the compromise faithfully---the
classic trusting-trust regress, answerable only by diversifying or
bootstrapping the toolchain itself. Reproducibility is thus one link
in a supply-chain integrity story that also needs source integrity
(version control with content-addressed history), dependency
integrity (pinned, hashed dependency manifests), and distribution
integrity (signed, logged releases). Each of these links has, since
roughly 2021, acquired a standardized and widely deployed
instantiation: the in-toto attestation framework and the SLSA
provenance levels \cite{intoto,slsa} give build provenance a signed,
machine-readable statement form, and the Sigstore
infrastructure \cite{sigstore}---keyless signing recorded in an
append-only transparency log---now backs release signing and
provenance at major open-source package registries. The survey
retains the generic vocabulary because the kernel adopts the
underlying ideas rather than these frameworks, but the frameworks
are where current practice lives. The general supply-chain lesson
for any evidence system is the one the kernel's compatibility
profile takes as a premise: an execution result is only as
meaningful as the identity of everything that participated in
producing it, and each participant---source bytes, interpreter,
environment, dependencies---must be either content-addressed and
recorded, or honestly declared out of scope. Declaring the
out-of-scope set is not an admission of failure; it is what makes
the in-scope claims exact. The dishonest alternative is a receipt
that is silent about what it did not pin.

\section{Provenance: entities, activities, derivations}\label{ch02:prov}

\subsection{The PROV model}\label{ch02:prov-model}

Where hashing answers ``is this the same content?'' and
reproducibility answers ``does this artifact correspond to that
source?'', provenance answers ``where did this come from, and
through what?'' The W3C PROV ontology \cite{prov}, which expresses
the PROV data model in OWL, standardizes the
vocabulary. Its core is a typed graph over three node kinds.
\emph{Entities} are things ``with some fixed aspects'' \cite{prov}---documents, datasets,
files, a particular version of a particular artifact.
\emph{Activities} are occurrents---things that, in the standard's
words, ``occur over a period of time and act upon or with
entities'' \cite{prov}. \emph{Agents} are things that ``bear some form
of responsibility'' \cite{prov}---people, organizations, or software to which activity is
attributed. The core edge kinds connect them: an activity
\emph{used} an entity; an entity \emph{was generated by} an
activity; an entity \emph{was derived from} another entity; an
entity \emph{was attributed to} an agent; an activity \emph{was
associated with} an agent. Derivation is the load-bearing relation
for evidence systems: it asserts that one entity's existence and
content depend on another's, and chains of derivation compose into
an auditable ancestry for any artifact.

Two modeling disciplines in PROV repay attention. First,
\emph{entities are versions, not names}: the model's ``fixed
aspects'' requirement means that a mutable thing is represented as a
sequence of entities, one per state---which harmonizes exactly with
content addressing, where each state has its own digest and hence
its own natural entity identifier. Second, \emph{activities are
first-class}: the run that produced an artifact is itself a node,
with its own attributes (start and end, parameters, environment),
not an annotation squeezed into the artifact. An evidence record
that identifies its generating activity---which suite, which
profile, which interpreter, which date---is a PROV generation edge
in typed-record form, and the kernel's execution receipts (later
chapters) are readable as exactly that.

\subsection{Granularity, bundles, and anti-patterns}\label{ch02:prov-granularity}

Applying the model well is mostly a matter of granularity decisions,
and the standard guidance can be compressed into three choices and
three anti-patterns. The choices: \emph{entity grain} (is the entity
the file, the record, or the field? finer grains support sharper
invalidation when an ancestor changes, at higher recording cost);
\emph{activity grain} (one activity per pipeline run, per stage, or
per operation? the right grain is the one at which parameters and
environment are actually uniform, since an activity's attributes
should be true of all of it); and \emph{assertion bundling} (PROV
allows provenance about provenance---who asserted this graph,
when---which is the hook by which provenance records themselves
enter the integrity-and-authority stack rather than floating above
it).

The anti-patterns are the ones audit experience keeps rediscovering.
\emph{Provenance by narrative}: ancestry recorded as prose in a
report rather than as typed edges over identified entities---
unqueryable, unverifiable, and quietly divergent from the artifacts
it describes. \emph{The mutable entity}: a provenance graph whose
nodes name mutable locations (``the config file'') rather than fixed
states, so that the graph's truth decays as the locations change;
the remedy is the content-addressed entity identifier already
discussed. \emph{The absent activity}: derivation edges recorded
without the generating activity, so that \emph{what was run}---the
question every serious audit asks first---is unanswerable; systems
that generate receipts as a side effect of execution, rather than
reconstructing them afterward, avoid this by construction. Each
anti-pattern has appeared, in some form, in the wider research
program's history, which is why its evidence tooling records typed,
digest-anchored records generated at execution time; the background
point is that these are the textbook remedies, not innovations.

\subsection{What provenance is not}\label{ch02:prov-limits}

Provenance is descriptive, not normative. A complete, honest
provenance graph for a worthless artifact is perfectly possible; so
is a beautiful ancestry terminating in a poisoned root. Provenance
assertions are, moreover, themselves just data---anyone can claim
that an entity was derived from another---so provenance without
integrity (content-addressed endpoints) is unverifiable narrative,
and provenance with integrity but without authorization (authenticated
asserters) is verifiable narrative of unknown authority. The three
layers stack: integrity fixes \emph{what} the nodes are, provenance
fixes \emph{how} they relate, and authorization---the subject of the
next section---fixes \emph{who stands behind} the assertions. Systems
get into trouble precisely when a lower layer is dressed as a higher
one: a derivation chain read as an endorsement chain, or a generation
record read as a review record. Keeping the layers typed and separate
is the design discipline; the vocabulary of \cite{prov} exists so
that the separation has standard names.

\section{Integrity versus authorization}\label{ch02:integrity-auth}

\subsection{Two different questions}\label{ch02:twoquestions}

Security engineering has long insisted on separating two questions
that casual system design merges \cite{anderson}. \emph{Integrity}:
is this artifact the artifact it claims to be---unmodified,
matching its digest, consistent with its schema?
\emph{Authorization}: is the action this artifact licenses actually
permitted---has a principal with the authority to permit it done so?
The questions have different verifiers, different failure modes, and
different machinery. Integrity is verified by computation over
content: rehash, recheck, reparse. Authorization is verified by
reference to a \emph{trust root}: some principal or key whose
authority is established outside the system under discussion---by
legal arrangement, organizational decision, or physical control---
and from which chains of delegation descend. A message
authentication code or a bare hash provides integrity; only a
construction bound to an authenticated principal---a signature
verifiable against a key whose holder is known and whose authority
over this action is established---begins to provide authorization,
and even then only relative to the policy that gives that key its
meaning.

The distinction is sharpest in its failure cases, which security
engineering catalogs at length \cite{anderson}. A system that
checks integrity and concludes authorization is vulnerable to
\emph{laundering}: any principal able to produce
internally-consistent artifacts---which is every principal, since
consistency is computable---can mint permissions. The archetype is
the self-signed assertion consumed as if countersigned: a document
that hashes correctly, validates against its schema, and says
``approved,'' where every one of those three facts was produced by
the same unaccountable process. Conversely, a system that checks
authorization but not integrity is vulnerable to substitution: a
genuinely authorized decision applied to swapped content. Sound
designs therefore verify both, separately, and are explicit about
which mechanism carries which burden. Policy names who may do what;
mechanism enforces it; and the enforcement of a policy must never be
inferable from properties that any party can self-supply.

\subsection{The mechanism inventory}\label{ch02:auth-mechanisms}

This section inventories, still as background, the standard
mechanisms that occupy the space between bare integrity and full
authorization, because their differences are exactly the
distinctions the kernel's records must not blur \cite{anderson}.

A \emph{bare hash} binds nothing to nobody: anyone can compute it,
so it provides integrity only relative to a trusted copy of the
digest itself. A \emph{message authentication code} binds content to
possession of a shared secret: it proves origin-or-tampering to the
other holder of the secret, but---because both parties can compute
it---provides no third-party-verifiable attribution and no
non-repudiation. A \emph{digital signature} binds content to a
private key: any holder of the public key can verify, and only the
private-key holder can sign, which yields attribution and
non-repudiation \emph{relative to the binding between key and
principal}. That binding is itself the next problem up: a
\emph{certificate} is a signed assertion by one key about the
identity or attributes of another, chains of certificates reduce
many bindings to few, and the chains terminate in roots distributed
out of band and defended by procedure---key ceremonies, hardware
custody, audit. Above all of this sits \emph{policy}: a valid
signature by an authentic key still authorizes nothing unless the
governing policy assigns that key authority over that action; the
verifier's real question is never ``is this signed?'' but ``is this
signed by a principal whom policy entitles to license this act?''

Laid out this way, the inventory exposes the category error at the
heart of the second failure class in the introduction: a self-hashed
artifact sits on the \emph{bottom} rung of this ladder---below the
MAC, since not even a shared secret is involved---while presenting
visual credentials (long hexadecimal strings, ``verified''
vocabulary) associated with the top. Nothing about the rungs is
subtle to a security engineer; the failure mode is that systems
without an explicit ladder let artifacts drift up it rhetorically.
Typed records that say which rung they occupy---this is a digest,
this is a signature, this key's authority derives from that
policy---are the standard countermeasure.

\subsection{Trust roots and the boundary of self-certification}\label{ch02:trustroots}

The deep point beneath the distinction is that authorization cannot
be bootstrapped from inside. Every authorization chain terminates in
a root that the verifying system takes on grounds external to
computation: a key ceremony, a legal identity, an organizational
appointment. A system with no such root can compute forever without
producing a single authorized bit; the most it can honestly produce
is tamper-evident, well-formed, thoroughly-described \emph{candidate}
material awaiting external review. This is not a defect to be
engineered around; it is the definition of the boundary. The honest
design consequence, where no trust root exists, is to make the
absence structural: represent authorization as a type the system
cannot inhabit, rather than as a flag the system could set. The
kernel's constant-false authorization stance, described in the
introduction and developed in the runtime chapter, is this standard
consequence applied without exception; the background point here is
that it is an application of textbook doctrine \cite{anderson}, not
an invention.

\section{Threat modeling}\label{ch02:threats}

\subsection{Method}\label{ch02:threat-method}

Threat modeling is the discipline of systematically asking what can
go wrong, before and during design rather than after deployment
\cite{shostack}. The method's canonical loop has four questions.
\emph{What are we building?}---answered with a model of the system:
its components, data flows, and, critically, its trust boundaries,
the surfaces where data or control crosses between principals of
different privilege or trustworthiness. \emph{What can go wrong?}---
answered by structured enumeration over the model. \emph{What are we
going to do about it?}---answered by mitigations: eliminate the
threatened element, add a control, or explicitly accept and document
the risk. \emph{Did we do a decent job?}---answered by validation:
checking the model against the built system and the mitigations
against the threats.

The best-known enumeration structure is STRIDE, a mnemonic over six
threat categories: \emph{Spoofing} (claiming an identity one does
not have), \emph{Tampering} (modifying data or code),
\emph{Repudiation} (denying an action absent evidence),
\emph{Information disclosure} (exposing data to those not entitled
to it), \emph{Denial of service} (degrading availability), and
\emph{Elevation of privilege} (gaining capability beyond one's
authorization) \cite{shostack}. Applied per element---each process,
data store, and data flow in the model is examined against each
applicable category---STRIDE converts the open question ``what can
go wrong?'' into a checklist with known blind spots rather than an
unbounded brainstorm. Its output is a set of \emph{threat
statements}: concrete, reviewable sentences of the form ``an actor
with access X performs Y, causing Z,'' each paired with a
mitigation, an acceptance, or an open item.

\subsection{Trust boundaries and model hygiene}\label{ch02:threat-boundaries}

The modeling step deserves its own comment, because the quality of
everything downstream is bounded by it. The standard modeling
artifact is the data-flow diagram: external entities, processes,
data stores, and the flows between them, with trust boundaries drawn
wherever a flow crosses between principals that should not fully
trust each other---user to service, service to database, parent
process to spawned worker, checked-out repository to executing
toolchain. Threat enumeration then concentrates on the crossings,
because that is where assumptions change hands: data validated on
one side is raw on the other; an identity authenticated in one zone
is a plain string in the next. Two hygiene rules from practice
\cite{shostack} transfer directly to evidence systems. First,
\emph{model the system you have, not the one you intend}: a diagram
omitting the debug endpoint, the cache, or the ``temporary'' shared
mutable status object will produce a threat list omitting their
threats, and audits repeatedly find that the omitted element is
where the incident lived. Second, \emph{boundaries are claims to be
enforced, not facts to be assumed}: drawing a line between worker
and parent means something only if a mechanism (a process boundary,
a sanitized channel, a validation gate) holds it, and every drawn
boundary should be traceable to its mechanism. A complementary
enumeration structure, attack trees---root goal decomposed into
conjunctive and disjunctive subgoals---is often used alongside
STRIDE when a single high-value failure (``an unauthorized claim is
consumed as authorized'') needs its paths enumerated exhaustively;
the two structures answer different questions (coverage per element
versus depth per goal) and mature models use both.

\subsection{Threats as review artifacts}\label{ch02:threat-artifacts}

Two features of the method matter for how this monograph is written.
First, threat statements are \emph{design inputs}, not post-hoc
documentation: a mechanism justified by a threat statement can be
reviewed for whether it actually blocks the stated threat, which is
a much sharper question than whether it is ``secure.'' Second, the
method treats \emph{explicit risk acceptance} as a first-class
outcome: a threat that is out of scope is written down as out of
scope, with the reasoning, rather than silently omitted---the
threat-model analog of the supply-chain rule that unpinned inputs
be declared. The later chapters of this monograph follow both
practices: each kernel mechanism is presented with the named threats
it answers (drawn from the program's actually-observed failure
classes, which serve as the threat catalog's empirical core), and
the non-goals---non-sandbox, non-hermetic, non-authoritative---are
presented as accepted, documented boundaries in exactly the sense
of \cite{shostack}. Where a STRIDE category maps cleanly onto a
kernel concern, the mapping is the expected one: spoofing maps to
identity and equality attacks, tampering to digest and registry
substitution, repudiation to receipt exactness, elevation of
privilege to claim promotion. The mapping is noted here once so
that later chapters can use the program's own vocabulary without
re-deriving it.

\subsection{Validation: closing the loop}\label{ch02:threat-validation}

The fourth question of the loop---did we do a decent job?---is the
one most often skipped, and its standard content is worth a
paragraph because the monograph's evidence chapters are, in effect,
one project's answer to it \cite{shostack}. Validation has two
halves. Model validation checks that the diagram still matches the
system: every component present, every flow drawn, every boundary's
enforcing mechanism still in place---a check that must recur,
since systems drift and threat models rot quietly. Mitigation
validation checks that each threat statement's mitigation actually
engages: for a fail-closed rule, that a test presents the
threatening input and observes the refusal; for an
unrepresentability claim, that review of the types finds no
constructor path; for an accepted risk, that the acceptance is
still visible in the artifacts consumers see (the false capability
field is present in the receipt, not merely in the design notes).
The discipline of pairing each threat with its validation evidence
---and of distinguishing ``mitigated by design,'' ``mitigated and
tested,'' and ``accepted and documented'' as three different
statuses---is what turns a threat model from a design-time essay
into a maintained contract. Where later chapters tabulate the
kernel's threats against its mechanisms and its receipts, they are
following this template, and the same template says exactly what
is missing when a validation cell is empty: an open obligation,
to be listed as such.

\section{Process isolation versus sandboxing}\label{ch02:isolation}

\subsection{A ladder of containment}\label{ch02:ladder}

``The code ran in isolation'' is among the most overloaded sentences
in systems engineering, and unpacking it matters wherever execution
receipts will be consumed as evidence. It is clarifying to arrange
the common meanings as a ladder, each rung adding a boundary the
previous rung lacks.

\emph{Rung 1: in-process hygiene.} The code runs in the same
process, with discipline: fresh module namespaces, controlled
globals. The weakest rung: it shares the entire address space,
interpreter state, and permissions of its host, and any misbehavior
is total.

\emph{Rung 2: process isolation.} The code runs in a separate
operating-system process, with a fresh address space and explicit
communication channels. For interpreted languages the rung includes
interpreter hygiene: launching the fresh interpreter with flags
that suppress inherited path injection, site customization, and
stale bytecode, and with a sanitized environment rather than the
parent's. Process isolation protects the \emph{parent's} state from
the child and makes the child's inputs enumerable. It does
\emph{not} constrain what the child may do to the world: the child
holds the same user privileges as its parent---files, network,
devices---and shares the host's native libraries and kernel.

\emph{Rung 3: OS sandboxing.} The child's privileges themselves are
reduced: system-call filtering, namespace and container isolation of
filesystem and network views, mandatory access control, resource
limits. A sandboxed child that misbehaves is constrained by the
operating system, not merely observed by its parent. Sandboxing is a
policy mechanism and is only as strong as the policy and the kernel
attack surface beneath it.

\emph{Rung 4: virtualization and hardware isolation.} The boundary
moves below the operating system: virtual machines, and further
down, separate physical hardware. Stronger again, with its own
residual channels (hypervisor bugs, side channels).

\emph{Orthogonal: hermeticity.} Cutting across the ladder is the
question of Section~\ref{ch02:repro}: is everything that
participates in the execution content-addressed? A hermetic
execution pins and records the identity of the interpreter, native
libraries, and OS surface, so that the receipt names its entire
closure. Hermeticity is orthogonal to containment---a strongly
sandboxed execution can be non-hermetic (unpinned toolchain inside
a strong box) and a weakly contained one can be partially hermetic
(pinned interpreter bytes, shared host).

\subsection{Interpreter-level isolation in practice}\label{ch02:interp}

Because the kernel's compatibility profile lives at rung 2, the
interpreter-hygiene details of that rung are worth spelling out in
their standard, non-project form; the discussion is phrased for
Python, whose mechanisms are representative. A child interpreter
launched naively inherits far more from its parent than the
process boundary suggests: environment variables that inject code
paths or site customizations; the working directory and script
directory prepended to the module search path, so that a file
named like a standard library module shadows it; user and system
site-packages directories, whose contents vary by machine and
time; startup files designated by environment variables; and stale
compiled bytecode caches that can desynchronize from their
sources. Each inheritance channel is a way for the execution's
behavior to depend on unrecorded host state---precisely the
supply-chain sin of Section~\ref{ch02:repro}---and each has a
standard closure: isolated mode (\texttt{-I}) removes the
script-directory and user-site injection and ignores the
environment variables that alter interpreter behavior; suppressing
site initialization (\texttt{-S}) removes machine-dependent path
customization; disabling bytecode writing (\texttt{-B}) prevents
cache pollution; and an explicitly constructed, minimal
environment replaces the inherited one. What remains, and must be
either recorded or declared unpinned, is the interpreter binary
itself, its standard library, and every native library it loads---
the recording of which is content addressing again, and the honest
declaration of which (when not recorded) is the capability-honesty
practice below.

One further rung-2 practice appears wherever isolated workers share
a host: concurrency discipline. Two workers serializing through a
shared resource---a cache directory, a lock file, a device---need a
mutual-exclusion mechanism, and \emph{advisory} locks (locks that
constrain only cooperating processes that choose to check them) are
the common tool. The standard honesty requirements transfer: an
advisory lock's guarantees should be stated as bounded (acquisition
timeouts, polling intervals) and unclaimed properties---fairness,
protection against non-cooperating processes---should be stated as
unclaimed. A receipt that says ``serialized via advisory lock,
without fairness claims'' is a rung-2 receipt telling the truth
about rung 2.

\subsection{Why the rung must be stated}\label{ch02:rung-honesty}

The engineering point of the ladder is not that higher is always
better---each rung has real costs in complexity, portability, and
performance, and the right rung depends on the threat model
(Section~\ref{ch02:threats}). The evidentiary point, which is the
one this monograph needs, is that \emph{a receipt must state its
rung}. An execution receipt that says ``isolated'' without saying
which boundary held invites the consumer to assume the strongest
reading; over time, the strongest reading becomes the folklore, and
the folklore becomes a false premise in someone's safety argument.
This is claim promotion (in the introduction's sense) applied to
containment. The standard remedy is capability honesty: receipts
carry explicit fields for the boundaries they did \emph{not}
establish---sandboxing absent, hermeticity absent---as first-class
negative facts, machine-readable, so that no downstream consumer
can honestly overread the rung. The kernel's compatibility receipts,
which expose their non-sandbox and non-hermetic status as explicit
false capability fields, instantiate this remedy; the remedy itself
is simply Section~\ref{ch02:threat-artifacts}'s documented-acceptance
practice transplanted to execution evidence.

\section{Synthesis: from background to kernel}\label{ch02:synthesis}

The seven bodies of material surveyed here compose into a single
picture, assembled here once before the mechanism
chapters descend into detail.

Hashing and content addressing (Section~\ref{ch02:hashing}) supply
the primitive: bytes can be named by content, and equal names mean
equal bytes. Canonical serialization (Section~\ref{ch02:canonical})
makes the primitive apply to structured values: with a total,
deterministic, injective encoding on a declared domain, equal
digests mean equal values and distinct values mean distinct
digests---provided the domain, number policy, string policy, type
tagging, and bounds are all fixed, which is why the strict
reject-and-tag family exists alongside interoperability-oriented
designs like JCS \cite{jcs}. Reproducible-build discipline
(Section~\ref{ch02:repro}) extends the same determinism from values
to executions: pin or record every participant, and declare what is
not pinned. Provenance (Section~\ref{ch02:prov}) supplies the
vocabulary for relating the resulting artifacts---entities generated
by activities, derived from entities, attributed to agents---without
ever confusing ancestry with endorsement. The
integrity-versus-authorization doctrine
(Section~\ref{ch02:integrity-auth}) then draws the line that the
whole stack must respect: everything so far is computable from
content and therefore self-suppliable, so none of it, however
elaborate, amounts to permission; permission requires a trust root
that computation cannot conjure. Threat modeling
(Section~\ref{ch02:threats}) provides the method by which a design
argues, mechanism by named threat, that its boundaries hold and
that its accepted risks are stated rather than hidden. And the
isolation ladder (Section~\ref{ch02:isolation}) applies the same
honesty to execution: say which boundary held, and carry the
negative capabilities in the receipt.

Table~\ref{ch02:tab-map} records the correspondence between this
background and the kernel mechanisms that later chapters develop.
The right-hand column is a map of forward pointers, not a claim
list; in every row, the idea is standard and the kernel's role is
adoption and composition under a fail-closed stance.

\begin{table}[htbp]
\centering
\caption{Background concepts and the kernel mechanisms that adopt
them.}
\label{ch02:tab-map}
\begin{tabular}{@{}ll@{}}
\toprule
Background concept & Kernel adoption (later chapters) \\
\midrule
Content addressing & definition digests; hash-bound artifacts \\
Canonical encoding (strict) & evidence encoder; tagged decimals \\
Injectivity & injective reference rendering \\
Reproducibility discipline & compiler-byte pinning; replay rule \\
Provenance (generation/use) & execution receipts; registry digests \\
Integrity vs.\ authorization & constant-false authorization \\
Trust roots & stated future-authorization constraints \\
Threat modeling & per-mechanism threat statements \\
Isolation ladder honesty & false capability fields in receipts \\
\bottomrule
\end{tabular}
\end{table}

\subsection{On the choice of references}\label{ch02:refs-rationale}

A brief note on why these five citations, and not fifty, anchor
the survey. Each was chosen as the consensus entry point to its
literature rather than as a research frontier. RFC~8785 \cite{jcs}
is the operative standard for JSON canonicalization---short,
normative, and implementable from its text, which makes it the
right fixed point for comparing stricter designs against.
The reproducible-builds literature is summarized through the
practitioner survey of Lamb and Zacchiroli \cite{reprobuilds},
which consolidates a decade of project experience---the failure
catalog, the rebuild-and-compare verification model, and the
deployment history---in citable form; for the adjacent supply-chain
attestation practice, the operative standards are the in-toto
attestation framework, the SLSA provenance levels, and the Sigstore
signing infrastructure \cite{intoto,slsa,sigstore}, named in
Section~\ref{ch02:repro-limits}. The W3C PROV ontology
\cite{prov} is the standardized vocabulary for provenance, and
citing the standard rather than the surrounding research keeps the
chapter's usage checkable against a fixed document. Anderson's
\emph{Security Engineering} \cite{anderson} is the standard
synthesis for the integrity/authentication/authorization
architecture this chapter leans on; the doctrine predates the
book, but the book is where it is assembled with the failure-case
evidence that gives it teeth. Shostack \cite{shostack} plays the
same role for threat modeling: the four-question frame and the
per-element STRIDE practice as taught to working engineers.
Topics for which the chapter cites nothing by name---hash-function
internals, XML canonicalization history, container isolation
mechanics---are treated at the level of common professional
knowledge, and readers wanting depth should consult the standard
texts of those areas; the chapter's obligations are to represent
the consensus accurately and to claim none of it, and both are
easier to audit when the reference list is short enough to
actually read.

\subsection{Composition pitfalls}\label{ch02:pitfalls}

A final piece of background: the ways in which correct components
of this stack compose incorrectly, since several kernel rules exist
to block specific compositions rather than specific components.
Four pitfalls recur in the literature and in audit experience.

\emph{The normalization sandwich.} Content is canonicalized, hashed,
then re-serialized by a downstream component with different rules,
and the re-serialization is later re-hashed for comparison. Both
serializers are individually deterministic; the composition is not.
Remedy: hash once at a declared boundary and transport the hashed
bytes, or re-canonicalize with the identical, versioned encoder---
never with a different one.

\emph{The identity switch.} One component compares raw digests,
another compares normalized digests (line endings, Unicode form,
case), and records flow between them. Each comparison is locally
sound; jointly they partition the artifact space differently, and
an artifact pair can be ``equal'' in one component and ``different''
in the next---an inconsistency an adversary can steer. Remedy: carry
each identity as its own named field, declare which field each
comparison uses, and forbid substitution---the two-identity
discipline noted in Section~\ref{ch02:contentaddr}.

\emph{The authority creep.} An integrity artifact is consumed by a
chain of components, each of which adds a small interpretive step,
until the last consumer treats it as authorization---the promotion
gradient of the introduction, restated as a composition failure.
Remedy: type the rungs (Section~\ref{ch02:auth-mechanisms}) so that
each component's input and output declare their rung, and any
rung-raising step requires the explicit machinery (a signature, a
policy check) rather than a paraphrase.

\emph{The boundary mirage.} A containment or validation boundary is
drawn in the design and assumed by downstream components, but the
mechanism enforcing it was weakened or removed---the sandbox that
became a process, the schema check that became a warning. Consumers
inherit a false premise. Remedy: capability honesty
(Section~\ref{ch02:rung-honesty}): receipts carry the enforced
boundary explicitly, including its negative fields, so that a
weakened boundary changes the receipts and the change is visible.

The common shape of all four is that composition erases context---
which encoder, which identity, which rung, which boundary---and the
common remedy is to make the context travel with the artifact as
typed fields that fail closed on mismatch. That sentence is, in
compressed form, the design brief the kernel sets itself; the
mechanisms of the following chapters are its elaboration.

\begin{boundary}[Non-novelty of this chapter]\label{ch02:bd-nonovel}
Nothing surveyed in this chapter is original to the CAS Core 8
project. Cryptographic hashing, content addressing, canonical
serialization and RFC~8785 \cite{jcs}, reproducible builds
\cite{reprobuilds}, the PROV provenance model \cite{prov}, the
integrity/authorization doctrine of security engineering
\cite{anderson}, threat-modeling method \cite{shostack}, and the
process-isolation/sandboxing distinction are community property,
cited here as standard background. Where later chapters describe
kernel mechanisms in this vocabulary, the reader should attribute
the vocabulary and the underlying techniques to this literature,
and to the kernel only the particular composition, the particular
domain restrictions, and the particular refusals it makes---each of
which remains, per the kernel's own boundary, an unreviewed Ring-4
research candidate and not a licensed or authoritative artifact.
\end{boundary}
